如上例所示, 为得到边界正则性结果 (也称为整体正则性结果), 需要假设边界具有某些光滑性. 为突出证明中最重要的技巧 (利用单位分解将区域分块, 并将含有部分边界的块映到半空间), 只对二阶标量方程给出证明, 且在证明中忽略低阶项.

\begin{Theorem}{整体正则性}
\label{thm:global-boundary-regularity}
假设 Theorem~\ref{thm:second-order-interior-regularity} 的假设成立, 并且 $\partial\Omega$ 属于 $C^2$ 类. 则 $u\in H^2(\Omega)$, 且
\begin{equation}\label{eq:global-boundary-regularity-h2-estimate}
\|u\|_{H^2(\Omega)}\leq C(\|u\|_{L^2(\Omega)}+\|f\|_{L^2(\Omega)}).
\end{equation}

\end{Theorem}

这个结果的证明较长且复杂, 因此将通过证明若干预备引理把它分开. 一个基本技巧是用单位分解将区域分块, 并将每一部分边界 ``展平\textquotesingle\textquotesingle. 正如第一个引理所示 (它实质上是边界已经平坦时主要结果的一个版本), 平坦边界使得差商技巧得以有效使用.

\begin{Lemma}
\label{lem:flat-boundary-local-regularity}
设 $R>0$, $\lambda\in(0,1)$, 定义
\begin{equation}\label{eq:flat-boundary-outer-half-ball}
D^+:=B_R(0)\cap\{\bm{x}\in\mathbb{R}^n\mid x_n>0\},
\end{equation}

\begin{equation}\label{eq:flat-boundary-inner-half-ball}
Q^+:=B_{\lambda R}(0)\cap\{\bm{x}\in\mathbb{R}^n\mid x_n>0\},
\end{equation}

\begin{equation}\label{eq:flat-boundary-source-isolated-empty-number}\end{equation}
设 $L$ 为形如下式的一致椭圆二阶算子\[L(\bm{x},D)u:=\sum_{i,j=1}^n\frac{\partial}{\partial x_i}a_{ij}(\bm{x})\frac{\partial u}{\partial x_j}+\sum_{i=1}^nb_i(\bm{x})\frac{\partial u}{\partial x_i}+c(\bm{x})u\]对应的双线性型为\[\begin{aligned}
B[v,u]&:=-\sum_{i,j=1}^n\int_{D^+}a_{ij}(\bm{x})u_{x_j}v_{x_i}\,\mathrm{d}\bm{x}+\sum_{i=1}^n\int_{D^+}b_i(\bm{x})u_{x_i}v\,\mathrm{d}\bm{x}\\
&\quad+\int_{D^+}c(\bm{x})uv\,\mathrm{d}\bm{x}.
\end{aligned}\]假设系数满足 $a_{ij}\in W^{1,\infty}(D^+)$, $b_i,c\in L^\infty(D^+)$, 且 $f\in L^2(D^+)$. 假设 $u\in H^1(D^+)$ 对所有 $v\in H_0^1(D^+)$ 满足变分方程
\begin{equation}\label{eq:flat-boundary-variational-equation}
B[v,u]=(v,f)
\end{equation}
并且在 $\{\bm{x}\in\mathbb{R}^n\mid x_n=0\}$ 上 $u\equiv0$ 在迹的意义下成立. 则 $u\in H^2(Q^+)$, 且存在依赖于 $R$ 的常数 $C$, 使得
\begin{equation}\label{eq:flat-boundary-local-h2-estimate}
\|u\|_{H^2(Q^+)}\leq C(\|f\|_{L^2(D^+)}+\|u\|_{L^2(D^+)}).
\end{equation}

\end{Lemma}
